% GENERATED FILE. Edit canonical lessons, then run npm run generate:books.
% canonical-source-hash: fnv1a64:ee77c6be9cba34da
% canonical-lessons: MW-C266-hajar, MW-C266-madad, MW-C266-meharbani, MW-C266-araj, MW-C266-asis

\chapter{Please Come In}
\label{ch:mw-please-come-in}
\hlchaptermodality{266}

\begin{chapteropening}
\textbf{By the end of this chapter:} \emph{I can greet a guest politely and ask for help.}

Complete the last lesson of chapter 266: I can greet a guest politely and ask for help.
\end{chapteropening}

\section[\texorpdfstring{\mw{हाजर} (hājar)}{हाजर (hājar)}]{\mw{हाजर} (hājar) — present; at your service}

\label{lesson:MW-C266-hajar}



\begin{quote}
Before the new one: say the Marwadi for brave, then the Marwadi for crazy.
\end{quote}

\subsection*{You'll want to know: \mw{हाजर}}
\textbf{\mw{हाजर}} — \emph{hājar} — \textquotedblleft{}present; at your service\textquotedblright{}.

\begin{quote}
\mw{हूं} \mw{हाजर} \mw{हूं} \mw{सा।}
\end{quote}

\begin{quote}
\emph{hū̃ hājar hū̃ sā.}
\end{quote}

\begin{quote}
I'm here, at your service.
\end{quote}

\begin{cousinweb}
\textbf{\mw{हाजर}} (\emph{hājar}) is usually traced, through Persian \emph{hāzir}, to Arabic \emph{ḥāḍir}, \textquotedblleft{}present\textquotedblright{}. With \textbf{\mw{सा}} it tells a host or an elder that you are here and ready to help.
\end{cousinweb}

\begin{grammarlens}[title={What you've built}]
One more word for asking and answering.
\end{grammarlens}

\subsection*{Your turn}
\begin{itemize}
\raggedright
  \item \emph{Say it:} \emph{hājar}
  \item \emph{Say it:} \emph{hājar}, once more
  \item \emph{Say it:} \emph{hū̃ hājar hū̃ sā.} — I'm here, at your service.
  \item \emph{Recall:} say the Marwadi for brave, then the Marwadi for crazy, then say \emph{hājar} again
\end{itemize}

\begin{tcolorbox}[breakable,colback=teal!4,colframe=teal!35!black,title={Before you move on}]
What does \mw{हाजर} mean? (\textquotedblleft{}Present; at your service\textquotedblright{}.) Say it once more.
\end{tcolorbox}
\section[\texorpdfstring{\mw{मदद} (madad)}{मदद (madad)}]{\mw{मदद} (madad) — help}

\label{lesson:MW-C266-madad}



\begin{quote}
Before the new one: say the Marwadi for crazy, then the Marwadi for at your service.
\end{quote}

\subsection*{You'll want to know: \mw{मदद}}
\textbf{\mw{मदद}} — \emph{madad} — \textquotedblleft{}help\textquotedblright{}.

\begin{quote}
\mw{आप} \mw{म्हारी} \mw{मदद} \mw{करो} \mw{सा।}
\end{quote}

\begin{quote}
\emph{āp mhārī madad karo sā.}
\end{quote}

\begin{quote}
Please help me.
\end{quote}

A feminine noun.

\begin{cousinweb}
\textbf{\mw{मदद}} (\emph{madad}) is usually traced to Persian \emph{madad}, \textquotedblleft{}help\textquotedblright{}. It is feminine: \textbf{\mw{म्हारी} \mw{मदद}}, my help.
\end{cousinweb}

\begin{grammarlens}[title={What you've built}]
One more word for asking and answering.
\end{grammarlens}

\subsection*{Your turn}
\begin{itemize}
\raggedright
  \item \emph{Say it:} \emph{madad}
  \item \emph{Say it:} \emph{madad}, once more
  \item \emph{Say it:} \emph{āp mhārī madad karo sā.} — Please help me.
  \item \emph{Recall:} say the Marwadi for crazy, then the Marwadi for at your service, then say \emph{madad} again
\end{itemize}

\begin{tcolorbox}[breakable,colback=teal!4,colframe=teal!35!black,title={Before you move on}]
What does \mw{मदद} mean? (\textquotedblleft{}Help\textquotedblright{}.) Say it once more.
\end{tcolorbox}
\section[\texorpdfstring{\mw{मेहरबानी} (meharbānī)}{मेहरबानी (meharbānī)}]{\mw{मेहरबानी} (meharbānī) — a kindness}

\label{lesson:MW-C266-meharbani}



\begin{quote}
Before the new one: say the Marwadi for at your service, then the Marwadi for help.
\end{quote}

\subsection*{You'll want to know: \mw{मेहरबानी}}
\textbf{\mw{मेहरबानी}} — \emph{meharbānī} — \textquotedblleft{}a kindness\textquotedblright{}.

\begin{quote}
\mw{घणी} \mw{मेहरबानी} \mw{सा}!
\end{quote}

\begin{quote}
\emph{ghaṇī meharbānī sā!}
\end{quote}

\begin{quote}
That's very kind of you!
\end{quote}

A feminine noun.

\begin{cousinweb}
\textbf{\mw{मेहरबानी}} (\emph{meharbānī}) is usually traced to Persian \emph{mihrbānī}, \textquotedblleft{}kindness\textquotedblright{}, a noun built on \emph{mihrbān}, \textquotedblleft{}kind\textquotedblright{}.
\end{cousinweb}

\begin{grammarlens}[title={What you've built}]
One more word for asking and answering.
\end{grammarlens}

\subsection*{Your turn}
\begin{itemize}
\raggedright
  \item \emph{Say it:} \emph{meharbānī}
  \item \emph{Say it:} \emph{meharbānī}, once more
  \item \emph{Say it:} \emph{ghaṇī meharbānī sā!} — That's very kind of you!
  \item \emph{Recall:} say the Marwadi for at your service, then the Marwadi for help, then say \emph{meharbānī} again
\end{itemize}

\begin{tcolorbox}[breakable,colback=teal!4,colframe=teal!35!black,title={Before you move on}]
What does \mw{मेहरबानी} mean? (\textquotedblleft{}A kindness\textquotedblright{}.) Say it once more.
\end{tcolorbox}
\section[\texorpdfstring{\mw{अरज} (araj)}{अरज (araj)}]{\mw{अरज} (araj) — a petition}

\label{lesson:MW-C266-araj}



\begin{quote}
Before the new one: say the Marwadi for help, then the Marwadi for a kindness.
\end{quote}

\subsection*{You'll want to know: \mw{अरज}}
\textbf{\mw{अरज}} — \emph{araj} — \textquotedblleft{}a petition\textquotedblright{}.

\begin{quote}
\mw{आप} \mw{म्हारी} \mw{अरज} \mw{सुणो} \mw{सा।}
\end{quote}

\begin{quote}
\emph{āp mhārī araj suṇo sā.}
\end{quote}

\begin{quote}
Please hear my petition.
\end{quote}

A feminine noun.

\begin{cousinweb}
\textbf{\mw{अरज}} (\emph{araj}) is usually traced, through Persian \emph{ʿarz}, to Arabic \emph{ʿarḍ}, \textquotedblleft{}presenting\textquotedblright{}: a request set before someone. It is feminine: \textbf{\mw{म्हारी} \mw{अरज}}.
\end{cousinweb}

\begin{grammarlens}[title={What you've built}]
One more word for asking and answering.
\end{grammarlens}

\subsection*{Your turn}
\begin{itemize}
\raggedright
  \item \emph{Say it:} \emph{araj}
  \item \emph{Say it:} \emph{araj}, once more
  \item \emph{Say it:} \emph{āp mhārī araj suṇo sā.} — Please hear my petition.
  \item \emph{Recall:} say the Marwadi for help, then the Marwadi for a kindness, then say \emph{araj} again
\end{itemize}

\begin{tcolorbox}[breakable,colback=teal!4,colframe=teal!35!black,title={Before you move on}]
What does \mw{अरज} mean? (\textquotedblleft{}A petition\textquotedblright{}.) Say it once more.
\end{tcolorbox}
\section[\texorpdfstring{\mw{आसीस} (āsīs)}{आसीस (āsīs)}]{\mw{आसीस} (āsīs) — a blessing}

\label{lesson:MW-C266-asis}



\begin{quote}
Before the new one: say the Marwadi for a kindness, then the Marwadi for a petition.
\end{quote}

\subsection*{You'll want to know: \mw{आसीस}}
\textbf{\mw{आसीस}} — \emph{āsīs} — \textquotedblleft{}a blessing\textquotedblright{}.

\begin{quote}
\mw{दादी} \mw{आसीस} \mw{देवै} \mw{है।}
\end{quote}

\begin{quote}
\emph{dādī āsīs devai hai.}
\end{quote}

\begin{quote}
Grandmother gives her blessing.
\end{quote}

A feminine noun.

\begin{cousinweb}
\textbf{\mw{आसीस}} (\emph{āsīs}) goes back to Sanskrit \emph{āśis}, \textquotedblleft{}a wish, a prayer; a blessing\textquotedblright{} (CDIAL 1457; Monier-Williams). CDIAL gives no Marwari form in this entry, and no other modern form either: the Marwadi word keeps the Sanskrit shape, with \mw{स} for \mw{श}. It is feminine, and an elder gives it: \textbf{\mw{आसीस} \mw{देवै} \mw{है}}.
\end{cousinweb}

\begin{grammarlens}[title={What you've built}]
One more word for asking and answering.
\end{grammarlens}

\subsection*{Your turn}
\begin{itemize}
\raggedright
  \item \emph{Say it:} \emph{āsīs}
  \item \emph{Say it:} \emph{āsīs}, once more
  \item \emph{Say it:} \emph{dādī āsīs devai hai.} — Grandmother gives her blessing.
  \item \emph{Recall:} say the Marwadi for a kindness, then the Marwadi for a petition, then say \emph{āsīs} again
\end{itemize}

\begin{tcolorbox}[breakable,colback=teal!4,colframe=teal!35!black,title={Before you move on}]
What does \mw{आसीस} mean? (\textquotedblleft{}A blessing\textquotedblright{}.) Say it once more.
\end{tcolorbox}
